\chapter{Introduction}
\label{ch:introduction}

\section{Motivation}
Mineral processing circuits such as grinding mills are multivariable, strongly
coupled and subject to frequent disturbances \cite{pomerleau2000}. They are
typically run with decentralized PID loops at the base layer and a supervisory
layer on top that coordinates setpoints and handles abnormal situations.
Learned controllers such as TD-MPC \cite{hansen2022tdmpc,hansen2024tdmpc2}
can perform well in this role, but their behaviour is encoded in neural
networks, which makes it hard for engineers to inspect, verify and certify
what the controller will do in situations it has not seen.
\TODO{Describe the industrial setting at Weir: the process, the current TD-MPC controller and why interpretability and auditability matter there.}

Recent work has shown that large language models (LLMs) can act as the
mutation operator in an evolutionary search over \emph{programs}, with an
automated evaluator as the only judge of quality
\cite{romeraparedes2024funsearch,novikov2025alphaevolve,ma2024eureka,weng2026lbg}.
The result of such a search is ordinary source code that can be read,
reviewed and tested like any other code. This thesis asks whether that idea
can produce supervisory control logic of practical quality.

\section{Research questions}
\begin{enumerate}
\item Can an LLM, used offline as a heuristic learner, generate deterministic,
  security-checked supervisory control code that performs close to a learned
  controller such as TD-MPC while remaining interpretable?
\item What information does the LLM need about the process, the evaluation and
  earlier attempts to make steady progress?
\item How must the evaluation be designed so that the search improves the
  intended behaviour rather than exploiting the score?
\item How much LLM reasoning is needed, and what does it cost?
\end{enumerate}
\TODO{Refine the research questions together with the supervisors.}

\section{Approach and contributions}
The work uses a three-layer architecture (Chapter~\ref{ch:method}): fixed PID
loops at the bottom, a deterministic supervisor function in the middle, and an
offline LLM-driven learner at the top that rewrites the supervisor between
runs. The approach is developed and evaluated on three toy processes of
increasing difficulty (Chapter~\ref{ch:toy}) before it is applied to the
industrial process. The contributions so far are:
\begin{itemize}
\item an implementation of the three-layer architecture with a security
  sandbox for generated code and a scoring scheme with safeguards against
  specification gaming;
\item results on a single tank, a two-tank cascade and the quadruple-tank
  process, including a comparison with a perfect-detection oracle;
\item an ablation of the LLM's reasoning effort on the quadruple-tank process.
\end{itemize}
\TODO{Add the industrial case study and the comparison with TD-MPC.}

\section{Outline}
Chapter~\ref{ch:background} gives background on supervisory process control,
fault detection and LLM-guided program search. Chapter~\ref{ch:method}
describes the architecture, the security sandbox, the scoring and the training
loop. Chapter~\ref{ch:toy} presents the three toy examples and their results.
Chapter~\ref{ch:discussion} discusses the findings and their limitations, and
Chapter~\ref{ch:conclusions} concludes.
